\documentclass[11pt,a4paper]{article}

\def\courseTitle{Industrial Security (Master)}
\def\labNumber{07}
\def\labTitle{Safety / Security Co-Engineering on a SIL-rated SIS -- Solutions}
\def\labSection{07}
\def\labTime{4 hours, group of 4--6}

% =====================================================================
% Shared preamble for lab sheets.
%
% Caller must \documentclass{article} first, then define:
%   \def\courseTitle{...}
%   \def\labNumber{NN}
%   \def\labTitle{...}
%   \def\labTime{e.g. "30--45 min"}
%   \def\labSection{e.g. "02"}
% Then % =====================================================================
% Shared preamble for lab sheets.
%
% Caller must \documentclass{article} first, then define:
%   \def\courseTitle{...}
%   \def\labNumber{NN}
%   \def\labTitle{...}
%   \def\labTime{e.g. "30--45 min"}
%   \def\labSection{e.g. "02"}
% Then % =====================================================================
% Shared preamble for lab sheets.
%
% Caller must \documentclass{article} first, then define:
%   \def\courseTitle{...}
%   \def\labNumber{NN}
%   \def\labTitle{...}
%   \def\labTime{e.g. "30--45 min"}
%   \def\labSection{e.g. "02"}
% Then \input{../../template/lab-preamble.tex}
% =====================================================================

\usepackage[utf8]{inputenc}
\usepackage[T1]{fontenc}
\usepackage[english]{babel}
\usepackage[a4paper, margin=2cm]{geometry}
\usepackage{graphicx}
\usepackage{listings}
\usepackage{xcolor}
\usepackage{colortbl}
\usepackage{tikz}
\usepackage{amsmath}
\usepackage{amssymb}
\usepackage{booktabs}
\usepackage{enumitem}
\usepackage{titlesec}
\usepackage{fancyhdr}
\usepackage{hyperref}
\usepackage{tcolorbox}
\tcbuselibrary{skins, breakable}
\usepackage{fontawesome5}

\input{../../../template/tikz-styles.tex}

\providecommand{\courseAuthor}{Industrial Security Lecture Series}
\providecommand{\courseInstitute}{Department of Computer Science}
\providecommand{\companyLogo}{logo}

\definecolor{slideAccent}{HTML}{00205B}
\definecolor{slideSecondary}{HTML}{00A0DC}

% --- Corporate branding overrides ------------------------------------
\IfFileExists{../../../branding.tex}{\input{../../../branding.tex}}{}

\colorlet{labAccent}{slideAccent}
\colorlet{labSec}{slideSecondary}
\definecolor{labCheck}{HTML}{2E7D32}
\definecolor{labWarn}{HTML}{B23A48}

\lstset{
  backgroundcolor=\color{black!4},
  basicstyle=\ttfamily\footnotesize,
  breaklines=true,
  frame=leftline,
  framerule=2pt,
  rulecolor=\color{labAccent},
  numbers=left,
  numberstyle=\tiny\color{black!50},
  showstringspaces=false,
  tabsize=2
}

\setlength{\tabcolsep}{4pt}

% Let TeX add a little extra inter-word stretch as a last resort before
% it reports an Overfull \hbox. Only kicks in on lines that would
% otherwise overflow (long \texttt paths, URLs, CIDRs); never loosens a
% line that already fits. Keeps prose/itemize within the margin.
\setlength{\emergencystretch}{3em}

\pagestyle{fancy}
\fancyhf{}
\renewcommand{\headrulewidth}{0.4pt}
\lhead{\small \courseTitle{} -- Lab \labNumber}
\rhead{\small Maps to Section \labSection}
\cfoot{\thepage}

\newtcolorbox{stepbox}[1]{
  enhanced, breakable, arc=2pt, boxrule=0pt,
  colback=labAccent!7, colframe=labAccent, leftrule=3pt,
  fonttitle=\bfseries\small, coltitle=white,
  colbacktitle=labAccent,
  title={\faTools\ Step #1}
}

\newtcolorbox{warnbox}{
  enhanced, breakable, arc=2pt, boxrule=0pt,
  colback=labWarn!7, colframe=labWarn, leftrule=4pt,
  fonttitle=\bfseries\small, coltitle=white, colbacktitle=labWarn,
  title={\faExclamationTriangle\ Caution}
}

\newtcolorbox{notebox}{
  enhanced, breakable, arc=2pt, boxrule=0pt,
  colback=labAccent!4, colframe=labAccent, leftrule=2pt,
  fonttitle=\bfseries\small, coltitle=white, colbacktitle=labAccent,
  title={\faInfoCircle\ Note}
}

\newtcolorbox{goalbox}{
  enhanced, breakable, arc=2pt, boxrule=0pt,
  colback=labCheck!7, colframe=labCheck, leftrule=3pt,
  fonttitle=\bfseries\small, coltitle=white, colbacktitle=labCheck,
  title={\faBullseye\ Goal}
}

\newtcolorbox{deliverbox}{
  enhanced, breakable, arc=2pt, boxrule=0pt,
  colback=labSec!10, colframe=labSec, leftrule=3pt,
  fonttitle=\bfseries\small, coltitle=white, colbacktitle=labSec,
  title={\faClipboardList\ Deliverable}
}

% --- Solution-sheet boxes (used only by lab-solutions.tex) ----------
\newtcolorbox{solutionbox}[1]{
  enhanced, breakable, arc=2pt, boxrule=0pt,
  colback=labCheck!7, colframe=labCheck, leftrule=3pt,
  fonttitle=\bfseries\small, coltitle=white, colbacktitle=labCheck,
  title={\faCheckCircle\ #1}
}

\newtcolorbox{rubricbox}{
  enhanced, breakable, arc=2pt, boxrule=0pt,
  colback=labSec!7, colframe=labSec, leftrule=3pt,
  fonttitle=\bfseries\small, coltitle=white, colbacktitle=labSec,
  title={\faBalanceScale\ Grading rubric}
}

\titleformat{\section}{\large\bfseries\color{labAccent}}{\thesection}{0.5em}{}

\newcommand{\labheader}{%
  \noindent
  \begin{tikzpicture}
    \fill[labAccent] (0,0) rectangle (\textwidth, -1.6cm);
    \node[anchor=north west, text=white, font=\bfseries\large] at (0.6cm, -0.4cm) {Lab \labNumber{} -- \labTitle};
    \node[anchor=south west, text=white, font=\small] at (0.6cm, -1.3cm) {\courseTitle{} \quad $\bullet$ \quad Maps to Section \labSection{} \quad $\bullet$ \quad \labTime};
    \node[anchor=north east, fill=labSec, text=white, font=\bfseries, inner sep=4pt, rounded corners=2pt]
      at ([xshift=-0.4cm, yshift=-0.4cm]\textwidth,0) {LAB \labNumber};
  \end{tikzpicture}
  \vspace{4mm}
}

% =====================================================================

\usepackage[utf8]{inputenc}
\usepackage[T1]{fontenc}
\usepackage[english]{babel}
\usepackage[a4paper, margin=2cm]{geometry}
\usepackage{graphicx}
\usepackage{listings}
\usepackage{xcolor}
\usepackage{colortbl}
\usepackage{tikz}
\usepackage{amsmath}
\usepackage{amssymb}
\usepackage{booktabs}
\usepackage{enumitem}
\usepackage{titlesec}
\usepackage{fancyhdr}
\usepackage{hyperref}
\usepackage{tcolorbox}
\tcbuselibrary{skins, breakable}
\usepackage{fontawesome5}

% =====================================================================
% Shared TikZ styles for the Industrial Security lecture series.
% Loaded by every slide deck and exercise sheet that uses TikZ.
% =====================================================================

\usetikzlibrary{
  arrows.meta,
  shapes,
  shapes.geometric,
  shapes.symbols,
  shapes.multipart,
  positioning,
  fit,
  calc,
  backgrounds,
  decorations.pathmorphing,
  decorations.pathreplacing,
  decorations.text,
  patterns,
  matrix,
  shadows
}

% --- colour palette --------------------------------------------------
\definecolor{isOT}{HTML}{1F4E79}     % OT (deep blue)
\definecolor{isIT}{HTML}{6F2DA8}     % IT (purple)
\definecolor{isSafety}{HTML}{C00000} % safety red
\definecolor{isNet}{HTML}{2E7D32}    % network green
\definecolor{isAttack}{HTML}{B23A48} % attack red
\definecolor{isAccent}{HTML}{E08E0B} % amber
\definecolor{isMute}{HTML}{6B6B6B}   % grey
\definecolor{isLight}{HTML}{EDF2F8}  % light background
\definecolor{isPLC}{HTML}{2C5282}
\definecolor{isHMI}{HTML}{2F855A}
\definecolor{isSCADA}{HTML}{6B46C1}

% --- generic node styles --------------------------------------------
\tikzset{
  isnode/.style={
    draw, rounded corners=2pt, minimum height=8mm, minimum width=18mm,
    font=\small, align=center, thick
  },
  isbox/.style={isnode, fill=isLight},
  plc/.style={isnode, fill=isPLC!15, draw=isPLC, thick},
  hmi/.style={isnode, fill=isHMI!15, draw=isHMI, thick},
  scada/.style={isnode, fill=isSCADA!15, draw=isSCADA, thick},
  sensor/.style={isnode, fill=isNet!10, draw=isNet, minimum height=6mm, minimum width=14mm, font=\scriptsize},
  actuator/.style={isnode, fill=isAccent!15, draw=isAccent, minimum height=6mm, minimum width=14mm, font=\scriptsize},
  firewall/.style={isnode, fill=isAttack!15, draw=isAttack, thick},
  server/.style={isnode, fill=isMute!15, draw=isMute!80, thick},
  switch/.style={isnode, fill=isLight, draw=black!60, minimum height=6mm, minimum width=14mm, font=\scriptsize},
  workstation/.style={isnode, fill=white, draw=isOT, thick},
  attacker/.style={isnode, fill=isAttack!20, draw=isAttack, thick, font=\small\bfseries},
  inetnode/.style={draw, ellipse, minimum width=24mm, minimum height=10mm, font=\scriptsize, fill=isLight, thick},
  process/.style={isnode, fill=isOT!15, draw=isOT, thick, minimum width=24mm},
  zone/.style={
    draw, rounded corners=4pt, thick, dashed, inner sep=4mm
  },
  conduit/.style={-{Stealth[length=2.5mm]}, thick, isNet!80!black},
  attack/.style={-{Stealth[length=2.5mm]}, thick, isAttack, dashed},
  flow/.style={-{Stealth[length=2.5mm]}, thick, isOT!70!black},
  netline/.style={thick, black!60},
  axisline/.style={-{Stealth[length=2mm]}, thick, black!70},
  tag/.style={font=\scriptsize\itshape, fill=white, inner sep=1pt},
  level/.style={
    draw, fill=isLight, minimum width=0.85\linewidth, minimum height=8mm,
    font=\small, anchor=west
  },
  brace/.style={decorate, decoration={brace, amplitude=4pt, raise=2pt}},
  legendbox/.style={draw, rounded corners=2pt, fill=isLight, inner sep=2mm, font=\scriptsize, align=left}
}

% --- handy macros ----------------------------------------------------
% \plcicon, \hmiicon etc as simple labelled boxes; useful inline.
\newcommand{\plcicon}[1][PLC]{\tikz[baseline=-0.6ex]{\node[plc, minimum width=10mm, minimum height=5mm, font=\scriptsize, inner sep=1pt]{#1};}}
\newcommand{\hmiicon}[1][HMI]{\tikz[baseline=-0.6ex]{\node[hmi, minimum width=10mm, minimum height=5mm, font=\scriptsize, inner sep=1pt]{#1};}}
\newcommand{\fwicon}[1][FW]{\tikz[baseline=-0.6ex]{\node[firewall, minimum width=10mm, minimum height=5mm, font=\scriptsize, inner sep=1pt]{#1};}}


\providecommand{\courseAuthor}{Industrial Security Lecture Series}
\providecommand{\courseInstitute}{Department of Computer Science}
\providecommand{\companyLogo}{logo}

\definecolor{slideAccent}{HTML}{00205B}
\definecolor{slideSecondary}{HTML}{00A0DC}

% --- Corporate branding overrides ------------------------------------
\IfFileExists{../../../branding.tex}{% =====================================================================
% Corporate / institutional branding -- single file to customise the
% look of the entire course (slides, notes, exercises, labs).
%
% HOW TO REBRAND IN UNDER A MINUTE:
%   1. Replace `branding/logo.pdf` with your own logo
%      (PDF preferred; PNG and JPG also work -- name it `logo.<ext>`).
%   2. (Optional) change the two colours below to your brand palette.
%   3. (Optional) override the author/institute lines below.
%   4. Re-run `make` at the repo root. That's it.
%
% Everything in this file has sensible defaults; you can delete a line
% to fall back to the built-in defaults, or comment it out.
%
% This file is loaded automatically by the slides, exercise, and lab
% preambles -- you never need to \input it manually.
% =====================================================================

% --- Logo --------------------------------------------------------------
% Path is resolved from each leaf-build directory; the default points at
% the shared `branding/` folder at the repo root. If you put your logo
% somewhere else, set the full or relative path here (without extension):
\renewcommand{\companyLogo}{../../../branding/logo}

% --- Author / institute (appears on the title page footer) ------------
\renewcommand{\courseAuthor}{}
\renewcommand{\courseInstitute}{}

% --- Brand colours ----------------------------------------------------
% slideAccent      = primary brand colour (title bands, accent rule)
% slideSecondary   = secondary brand colour (section badge, highlight)
% Both use HTML hex codes. Keep enough contrast against white text.
\definecolor{slideAccent}{HTML}{00205B}      % deep navy
\definecolor{slideSecondary}{HTML}{00A0DC}   % light blue accent
}{}

\colorlet{labAccent}{slideAccent}
\colorlet{labSec}{slideSecondary}
\definecolor{labCheck}{HTML}{2E7D32}
\definecolor{labWarn}{HTML}{B23A48}

\lstset{
  backgroundcolor=\color{black!4},
  basicstyle=\ttfamily\footnotesize,
  breaklines=true,
  frame=leftline,
  framerule=2pt,
  rulecolor=\color{labAccent},
  numbers=left,
  numberstyle=\tiny\color{black!50},
  showstringspaces=false,
  tabsize=2
}

\setlength{\tabcolsep}{4pt}

% Let TeX add a little extra inter-word stretch as a last resort before
% it reports an Overfull \hbox. Only kicks in on lines that would
% otherwise overflow (long \texttt paths, URLs, CIDRs); never loosens a
% line that already fits. Keeps prose/itemize within the margin.
\setlength{\emergencystretch}{3em}

\pagestyle{fancy}
\fancyhf{}
\renewcommand{\headrulewidth}{0.4pt}
\lhead{\small \courseTitle{} -- Lab \labNumber}
\rhead{\small Maps to Section \labSection}
\cfoot{\thepage}

\newtcolorbox{stepbox}[1]{
  enhanced, breakable, arc=2pt, boxrule=0pt,
  colback=labAccent!7, colframe=labAccent, leftrule=3pt,
  fonttitle=\bfseries\small, coltitle=white,
  colbacktitle=labAccent,
  title={\faTools\ Step #1}
}

\newtcolorbox{warnbox}{
  enhanced, breakable, arc=2pt, boxrule=0pt,
  colback=labWarn!7, colframe=labWarn, leftrule=4pt,
  fonttitle=\bfseries\small, coltitle=white, colbacktitle=labWarn,
  title={\faExclamationTriangle\ Caution}
}

\newtcolorbox{notebox}{
  enhanced, breakable, arc=2pt, boxrule=0pt,
  colback=labAccent!4, colframe=labAccent, leftrule=2pt,
  fonttitle=\bfseries\small, coltitle=white, colbacktitle=labAccent,
  title={\faInfoCircle\ Note}
}

\newtcolorbox{goalbox}{
  enhanced, breakable, arc=2pt, boxrule=0pt,
  colback=labCheck!7, colframe=labCheck, leftrule=3pt,
  fonttitle=\bfseries\small, coltitle=white, colbacktitle=labCheck,
  title={\faBullseye\ Goal}
}

\newtcolorbox{deliverbox}{
  enhanced, breakable, arc=2pt, boxrule=0pt,
  colback=labSec!10, colframe=labSec, leftrule=3pt,
  fonttitle=\bfseries\small, coltitle=white, colbacktitle=labSec,
  title={\faClipboardList\ Deliverable}
}

% --- Solution-sheet boxes (used only by lab-solutions.tex) ----------
\newtcolorbox{solutionbox}[1]{
  enhanced, breakable, arc=2pt, boxrule=0pt,
  colback=labCheck!7, colframe=labCheck, leftrule=3pt,
  fonttitle=\bfseries\small, coltitle=white, colbacktitle=labCheck,
  title={\faCheckCircle\ #1}
}

\newtcolorbox{rubricbox}{
  enhanced, breakable, arc=2pt, boxrule=0pt,
  colback=labSec!7, colframe=labSec, leftrule=3pt,
  fonttitle=\bfseries\small, coltitle=white, colbacktitle=labSec,
  title={\faBalanceScale\ Grading rubric}
}

\titleformat{\section}{\large\bfseries\color{labAccent}}{\thesection}{0.5em}{}

\newcommand{\labheader}{%
  \noindent
  \begin{tikzpicture}
    \fill[labAccent] (0,0) rectangle (\textwidth, -1.6cm);
    \node[anchor=north west, text=white, font=\bfseries\large] at (0.6cm, -0.4cm) {Lab \labNumber{} -- \labTitle};
    \node[anchor=south west, text=white, font=\small] at (0.6cm, -1.3cm) {\courseTitle{} \quad $\bullet$ \quad Maps to Section \labSection{} \quad $\bullet$ \quad \labTime};
    \node[anchor=north east, fill=labSec, text=white, font=\bfseries, inner sep=4pt, rounded corners=2pt]
      at ([xshift=-0.4cm, yshift=-0.4cm]\textwidth,0) {LAB \labNumber};
  \end{tikzpicture}
  \vspace{4mm}
}

% =====================================================================

\usepackage[utf8]{inputenc}
\usepackage[T1]{fontenc}
\usepackage[english]{babel}
\usepackage[a4paper, margin=2cm]{geometry}
\usepackage{graphicx}
\usepackage{listings}
\usepackage{xcolor}
\usepackage{colortbl}
\usepackage{tikz}
\usepackage{amsmath}
\usepackage{amssymb}
\usepackage{booktabs}
\usepackage{enumitem}
\usepackage{titlesec}
\usepackage{fancyhdr}
\usepackage{hyperref}
\usepackage{tcolorbox}
\tcbuselibrary{skins, breakable}
\usepackage{fontawesome5}

% =====================================================================
% Shared TikZ styles for the Industrial Security lecture series.
% Loaded by every slide deck and exercise sheet that uses TikZ.
% =====================================================================

\usetikzlibrary{
  arrows.meta,
  shapes,
  shapes.geometric,
  shapes.symbols,
  shapes.multipart,
  positioning,
  fit,
  calc,
  backgrounds,
  decorations.pathmorphing,
  decorations.pathreplacing,
  decorations.text,
  patterns,
  matrix,
  shadows
}

% --- colour palette --------------------------------------------------
\definecolor{isOT}{HTML}{1F4E79}     % OT (deep blue)
\definecolor{isIT}{HTML}{6F2DA8}     % IT (purple)
\definecolor{isSafety}{HTML}{C00000} % safety red
\definecolor{isNet}{HTML}{2E7D32}    % network green
\definecolor{isAttack}{HTML}{B23A48} % attack red
\definecolor{isAccent}{HTML}{E08E0B} % amber
\definecolor{isMute}{HTML}{6B6B6B}   % grey
\definecolor{isLight}{HTML}{EDF2F8}  % light background
\definecolor{isPLC}{HTML}{2C5282}
\definecolor{isHMI}{HTML}{2F855A}
\definecolor{isSCADA}{HTML}{6B46C1}

% --- generic node styles --------------------------------------------
\tikzset{
  isnode/.style={
    draw, rounded corners=2pt, minimum height=8mm, minimum width=18mm,
    font=\small, align=center, thick
  },
  isbox/.style={isnode, fill=isLight},
  plc/.style={isnode, fill=isPLC!15, draw=isPLC, thick},
  hmi/.style={isnode, fill=isHMI!15, draw=isHMI, thick},
  scada/.style={isnode, fill=isSCADA!15, draw=isSCADA, thick},
  sensor/.style={isnode, fill=isNet!10, draw=isNet, minimum height=6mm, minimum width=14mm, font=\scriptsize},
  actuator/.style={isnode, fill=isAccent!15, draw=isAccent, minimum height=6mm, minimum width=14mm, font=\scriptsize},
  firewall/.style={isnode, fill=isAttack!15, draw=isAttack, thick},
  server/.style={isnode, fill=isMute!15, draw=isMute!80, thick},
  switch/.style={isnode, fill=isLight, draw=black!60, minimum height=6mm, minimum width=14mm, font=\scriptsize},
  workstation/.style={isnode, fill=white, draw=isOT, thick},
  attacker/.style={isnode, fill=isAttack!20, draw=isAttack, thick, font=\small\bfseries},
  inetnode/.style={draw, ellipse, minimum width=24mm, minimum height=10mm, font=\scriptsize, fill=isLight, thick},
  process/.style={isnode, fill=isOT!15, draw=isOT, thick, minimum width=24mm},
  zone/.style={
    draw, rounded corners=4pt, thick, dashed, inner sep=4mm
  },
  conduit/.style={-{Stealth[length=2.5mm]}, thick, isNet!80!black},
  attack/.style={-{Stealth[length=2.5mm]}, thick, isAttack, dashed},
  flow/.style={-{Stealth[length=2.5mm]}, thick, isOT!70!black},
  netline/.style={thick, black!60},
  axisline/.style={-{Stealth[length=2mm]}, thick, black!70},
  tag/.style={font=\scriptsize\itshape, fill=white, inner sep=1pt},
  level/.style={
    draw, fill=isLight, minimum width=0.85\linewidth, minimum height=8mm,
    font=\small, anchor=west
  },
  brace/.style={decorate, decoration={brace, amplitude=4pt, raise=2pt}},
  legendbox/.style={draw, rounded corners=2pt, fill=isLight, inner sep=2mm, font=\scriptsize, align=left}
}

% --- handy macros ----------------------------------------------------
% \plcicon, \hmiicon etc as simple labelled boxes; useful inline.
\newcommand{\plcicon}[1][PLC]{\tikz[baseline=-0.6ex]{\node[plc, minimum width=10mm, minimum height=5mm, font=\scriptsize, inner sep=1pt]{#1};}}
\newcommand{\hmiicon}[1][HMI]{\tikz[baseline=-0.6ex]{\node[hmi, minimum width=10mm, minimum height=5mm, font=\scriptsize, inner sep=1pt]{#1};}}
\newcommand{\fwicon}[1][FW]{\tikz[baseline=-0.6ex]{\node[firewall, minimum width=10mm, minimum height=5mm, font=\scriptsize, inner sep=1pt]{#1};}}


\providecommand{\courseAuthor}{Industrial Security Lecture Series}
\providecommand{\courseInstitute}{Department of Computer Science}
\providecommand{\companyLogo}{logo}

\definecolor{slideAccent}{HTML}{00205B}
\definecolor{slideSecondary}{HTML}{00A0DC}

% --- Corporate branding overrides ------------------------------------
\IfFileExists{../../../branding.tex}{% =====================================================================
% Corporate / institutional branding -- single file to customise the
% look of the entire course (slides, notes, exercises, labs).
%
% HOW TO REBRAND IN UNDER A MINUTE:
%   1. Replace `branding/logo.pdf` with your own logo
%      (PDF preferred; PNG and JPG also work -- name it `logo.<ext>`).
%   2. (Optional) change the two colours below to your brand palette.
%   3. (Optional) override the author/institute lines below.
%   4. Re-run `make` at the repo root. That's it.
%
% Everything in this file has sensible defaults; you can delete a line
% to fall back to the built-in defaults, or comment it out.
%
% This file is loaded automatically by the slides, exercise, and lab
% preambles -- you never need to \input it manually.
% =====================================================================

% --- Logo --------------------------------------------------------------
% Path is resolved from each leaf-build directory; the default points at
% the shared `branding/` folder at the repo root. If you put your logo
% somewhere else, set the full or relative path here (without extension):
\renewcommand{\companyLogo}{../../../branding/logo}

% --- Author / institute (appears on the title page footer) ------------
\renewcommand{\courseAuthor}{}
\renewcommand{\courseInstitute}{}

% --- Brand colours ----------------------------------------------------
% slideAccent      = primary brand colour (title bands, accent rule)
% slideSecondary   = secondary brand colour (section badge, highlight)
% Both use HTML hex codes. Keep enough contrast against white text.
\definecolor{slideAccent}{HTML}{00205B}      % deep navy
\definecolor{slideSecondary}{HTML}{00A0DC}   % light blue accent
}{}

\colorlet{labAccent}{slideAccent}
\colorlet{labSec}{slideSecondary}
\definecolor{labCheck}{HTML}{2E7D32}
\definecolor{labWarn}{HTML}{B23A48}

\lstset{
  backgroundcolor=\color{black!4},
  basicstyle=\ttfamily\footnotesize,
  breaklines=true,
  frame=leftline,
  framerule=2pt,
  rulecolor=\color{labAccent},
  numbers=left,
  numberstyle=\tiny\color{black!50},
  showstringspaces=false,
  tabsize=2
}

\setlength{\tabcolsep}{4pt}

% Let TeX add a little extra inter-word stretch as a last resort before
% it reports an Overfull \hbox. Only kicks in on lines that would
% otherwise overflow (long \texttt paths, URLs, CIDRs); never loosens a
% line that already fits. Keeps prose/itemize within the margin.
\setlength{\emergencystretch}{3em}

\pagestyle{fancy}
\fancyhf{}
\renewcommand{\headrulewidth}{0.4pt}
\lhead{\small \courseTitle{} -- Lab \labNumber}
\rhead{\small Maps to Section \labSection}
\cfoot{\thepage}

\newtcolorbox{stepbox}[1]{
  enhanced, breakable, arc=2pt, boxrule=0pt,
  colback=labAccent!7, colframe=labAccent, leftrule=3pt,
  fonttitle=\bfseries\small, coltitle=white,
  colbacktitle=labAccent,
  title={\faTools\ Step #1}
}

\newtcolorbox{warnbox}{
  enhanced, breakable, arc=2pt, boxrule=0pt,
  colback=labWarn!7, colframe=labWarn, leftrule=4pt,
  fonttitle=\bfseries\small, coltitle=white, colbacktitle=labWarn,
  title={\faExclamationTriangle\ Caution}
}

\newtcolorbox{notebox}{
  enhanced, breakable, arc=2pt, boxrule=0pt,
  colback=labAccent!4, colframe=labAccent, leftrule=2pt,
  fonttitle=\bfseries\small, coltitle=white, colbacktitle=labAccent,
  title={\faInfoCircle\ Note}
}

\newtcolorbox{goalbox}{
  enhanced, breakable, arc=2pt, boxrule=0pt,
  colback=labCheck!7, colframe=labCheck, leftrule=3pt,
  fonttitle=\bfseries\small, coltitle=white, colbacktitle=labCheck,
  title={\faBullseye\ Goal}
}

\newtcolorbox{deliverbox}{
  enhanced, breakable, arc=2pt, boxrule=0pt,
  colback=labSec!10, colframe=labSec, leftrule=3pt,
  fonttitle=\bfseries\small, coltitle=white, colbacktitle=labSec,
  title={\faClipboardList\ Deliverable}
}

% --- Solution-sheet boxes (used only by lab-solutions.tex) ----------
\newtcolorbox{solutionbox}[1]{
  enhanced, breakable, arc=2pt, boxrule=0pt,
  colback=labCheck!7, colframe=labCheck, leftrule=3pt,
  fonttitle=\bfseries\small, coltitle=white, colbacktitle=labCheck,
  title={\faCheckCircle\ #1}
}

\newtcolorbox{rubricbox}{
  enhanced, breakable, arc=2pt, boxrule=0pt,
  colback=labSec!7, colframe=labSec, leftrule=3pt,
  fonttitle=\bfseries\small, coltitle=white, colbacktitle=labSec,
  title={\faBalanceScale\ Grading rubric}
}

\titleformat{\section}{\large\bfseries\color{labAccent}}{\thesection}{0.5em}{}

\newcommand{\labheader}{%
  \noindent
  \begin{tikzpicture}
    \fill[labAccent] (0,0) rectangle (\textwidth, -1.6cm);
    \node[anchor=north west, text=white, font=\bfseries\large] at (0.6cm, -0.4cm) {Lab \labNumber{} -- \labTitle};
    \node[anchor=south west, text=white, font=\small] at (0.6cm, -1.3cm) {\courseTitle{} \quad $\bullet$ \quad Maps to Section \labSection{} \quad $\bullet$ \quad \labTime};
    \node[anchor=north east, fill=labSec, text=white, font=\bfseries, inner sep=4pt, rounded corners=2pt]
      at ([xshift=-0.4cm, yshift=-0.4cm]\textwidth,0) {LAB \labNumber};
  \end{tikzpicture}
  \vspace{4mm}
}


\begin{document}
\labheader

\begin{goalbox}
Lecturer / TA reference for the HF-201 co-engineering workshop: a worked
Cyber-HAZOP row set, the reference $\mathrm{PFD}_{avg}$ / SIL calculation,
the LOPA-independence-defeat analysis with before/after numbers, model
non-interfering controls, the lifecycle mapping, a GSN sketch, common
student mistakes, and a grading rubric. Treat the numbers as the reference
answer; full marks for any defensible alternative with sound reasoning.
\end{goalbox}

\begin{solutionbox}{Steps 2--3 -- Worked Cyber-HAZOP rows (paired causes, phase tag)}
Accept any six credible rows; at least three cyber causes must target the SIS
itself. The set below is the reference standard. \emph{Phase} is the
TR84.00.09 lifecycle phase the cyber cause exploits.

\vspace{2pt}\footnotesize
\renewcommand{\arraystretch}{1.25}\setlength{\tabcolsep}{3pt}
\begin{tabular}{@{}p{10mm}p{14mm}p{30mm}p{30mm}p{28mm}p{14mm}@{}}
\toprule
\textbf{Guide} & \textbf{Param} & \textbf{Safety cause} & \textbf{Cyber cause} & \textbf{Safeguard / gap} & \textbf{Phase} \\
\midrule
More & Header press. & PT drift / blocked impulse line & Setpoint raised via BPCS HMI write & SIF-1 trip; gap: no write allow-list & Oper. \\
No   & Fuel isolation & Solenoid sticks open & Logic appended so DBB valves never command closed & 1oo1 FE; gap: unsigned download & Realis. \\
As~well~as & Flame status & Scanner fouling & 2oo3 vote spoofed to ``flame OK'' on SIS I/O & 2oo3 voting; gap: no I/O integrity check & Design \\
Other~than & SIS mode & Random CPU fault & Key-switch left in PROGRAM, logic overwritten & Key-switch interlock; gap: PROGRAM not alarmed & Oper. \\
Part~of & Safety net & PSU degradation & Flat LAN lets EWS reach both SIS and BPCS & Zone firewall absent; gap: no conduit segmentation & Design \\
More & Demand rate & Process upset & MOC bypassed, SIF disabled without re-validation & Change board; gap: no security sign-off on MOC & Modif. \\
\bottomrule
\end{tabular}
\normalsize

\par\vspace{2mm}
\textbf{Marking note.} The discriminating skill is the phase tag. Rows tagged
\emph{Operation} or \emph{Modification} are the high-value findings: they
defeat the safety function \emph{after} certification, so the SIL certificate
no longer evidences integrity. A worksheet with only \emph{Design}-phase
network gripes has missed the point.
\end{solutionbox}

\begin{solutionbox}{Step 4 -- Reference $\mathrm{PFD}_{avg}$ / SIL calculation (SIF-1)}
With $1\,\mathrm{FIT}=10^{-9}\,\mathrm{h}^{-1}$ and $T_I=8760\,\mathrm{h}$:

\par\vspace{1mm}
\textbf{Sensor (1oo2 PTs, $\lambda_{DU}=300\,\mathrm{FIT}=3.0\times10^{-7}$, $\beta=0.10$):}
\[
\mathrm{PFD}_S \approx \tfrac{1}{3}(\lambda_{DU}T_I)^2 + \beta\,\lambda_{DU}\tfrac{T_I}{2}
\]
$\lambda_{DU}T_I = 3.0\times10^{-7}\times8760 = 2.628\times10^{-3}$.
First term $=\tfrac13(2.628\times10^{-3})^2 = 2.30\times10^{-6}$;
common-cause term $=0.10\times3.0\times10^{-7}\times4380 = 1.314\times10^{-4}$.
So $\mathrm{PFD}_S \approx 1.34\times10^{-4}$ (the $\beta$ term dominates --
the teaching point of redundancy with common cause).

\par\vspace{1mm}
\textbf{Logic solver (TMR, certified):} $\mathrm{PFD}_L = 2.0\times10^{-4}$.

\par\vspace{1mm}
\textbf{Final element (1oo1, $\lambda_{DU}=600\,\mathrm{FIT}=6.0\times10^{-7}$):}
\[
\mathrm{PFD}_{FE} \approx \tfrac12\lambda_{DU}T_I = \tfrac12\times6.0\times10^{-7}\times8760 = 2.63\times10^{-3}.
\]

\par\vspace{1mm}
\textbf{Total:}
\[
\mathrm{PFD}_{avg}^{SIF\text{-}1} = 1.34\times10^{-4} + 2.0\times10^{-4} + 2.63\times10^{-3} \approx 2.96\times10^{-3}.
\]
This is in $[10^{-3},10^{-2})$ $\Rightarrow$ \textbf{SIL\,2 achieved}
(RRF $\approx 1/2.96\times10^{-3}\approx 338$). The 1oo1 final element
dominates -- correct engineering conclusion is that improving the SIF means
1oo2 valves or a shorter proof interval, not touching the logic solver.
Accept answers within a factor of $\sim$2 given rounding.
\end{solutionbox}

\begin{solutionbox}{Step 5 -- LOPA-independence defeat (the central finding)}
\textbf{Per-layer verdict after the SIS-EWS is compromised:}
\begin{itemize}[leftmargin=*, nosep]
  \item \textbf{SIS / SIF-1:} \emph{defeated.} Unsigned download + key-switch
        in PROGRAM let the attacker append logic. Credit $\to$ PFD $=1$.
  \item \textbf{BPCS alarm + operator:} \emph{defeated.} Flat LAN shared with
        the SIS means the same foothold reaches the BPCS; the attacker
        suppresses the high-pressure alarm. Credit $\to$ PFD $=1$.
  \item \textbf{Mechanical relief valve:} \emph{survives.} Purely mechanical,
        outside the cyber blast radius. Credit retained, PFD $\approx 10^{-2}$.
\end{itemize}

\par\vspace{1mm}
\textbf{Independence collapse:} the SIS and the BPCS-alarm layer were credited
as independent, but the \emph{shared engineering path + flat network} is a
single common cause that defeats both. LOPA's independence axiom is violated
the moment one foothold can reach both layers. The relief valve is the only
genuinely independent layer left.

\par\vspace{1mm}
\textbf{Before/after (illustrative, $f_{IE}=10^{-1}$/yr exothermic upset):}
\[
f_{TE}^{\text{before}} = 10^{-1}\times(3.0\times10^{-3})\times(10^{-1})\times(10^{-2}) \approx 3\times10^{-7}\,/\text{yr}
\]
\[
f_{TE}^{\text{after}} = 10^{-1}\times 1 \times 1 \times (10^{-2}) = 10^{-3}\,/\text{yr}
\]
A loss of $\sim$3--4 orders of magnitude of protection. The point students
must articulate: the danger is not the malware payload, it is that
\emph{counting correlated layers as independent} produced a SIL claim that was
never true against a cyber adversary.
\end{solutionbox}

\begin{solutionbox}{Step 6 -- TRITON/TRISIS walk-through + safe-state}
Model answer per stage:
\begin{itemize}[leftmargin=*, nosep]
  \item \textbf{Foothold $\to$ EWS:} dual-homed NIC bridges DMZ to safety net.
        Stopped by a data diode for the historian feed (no inbound path) and
        a zone firewall.
  \item \textbf{EWS $\to$ controller:} proprietary protocol + key in PROGRAM
        allow logic append. Detected by key-switch-in-RUN + alarm on PROGRAM,
        and by independent code-integrity (checksum) monitoring.
  \item \textbf{Payload:} TRITON appended a backdoor intending to let unsafe
        states pass while the SIS appeared healthy; the real incident tripped
        because a TMR memory-validation mismatch forced a safe fault state.
\end{itemize}
\textbf{Was the trip luck or design?} Both. Fail-safe TMR \emph{design} means a
detected internal inconsistency drives the safe state -- that is what tripped
the plant. But the attacker did not \emph{intend} a trip; the protective
behaviour was triggered by a payload defect, not by the SIS detecting the
attack. So fail-safe TMR buys you a safe state on \emph{detected} faults; it
does \emph{not} buy you protection against a clean payload that keeps the
controller internally consistent while disabling the protective response (the
``degraded but running'' attack). Safe-state implication: after a successful,
defect-free version of this payload, the SIS would \emph{not} reach safe state
on a real demand -- which is exactly the catastrophic outcome the attacker
sought. The trip was a near miss, not evidence the SIS was secure.
\end{solutionbox}

\begin{solutionbox}{Step 7 -- Model non-interfering controls (accepted vs.\ rejected)}
\textbf{Accepted -- each removes a threat without touching the safety path:}
\begin{itemize}[leftmargin=*, nosep]
  \item \textbf{Key-switch to RUN, dual-control custody + PROGRAM alarm.}
        Removes the post-cert logic-overwrite path. Cannot raise
        $\mathrm{PFD}_{avg}$ (it does not alter the SIF); a genuine emergency
        download is still possible under control. No re-validation needed.
  \item \textbf{Zone the SIS off the BPCS (IEC~62443-3-2 conduit + firewall).}
        Breaks the flat-LAN common cause from Step~5. Passive to the safety
        logic. No SIF change $\Rightarrow$ no re-validation.
  \item \textbf{Data diode for the historian feed; remove the dual-homed NIC.}
        Eliminates the inbound DMZ$\to$SIS path. One-way egress only; cannot
        affect safety output.
  \item \textbf{Signed / whitelisted logic downloads.} Closes the unsigned
        download gap. Is part of the engineering workflow, not the runtime
        safety loop; does not add latency in the process safety time.
  \item \textbf{Independent (out-of-band, passive) code-integrity monitoring.}
        Detects the appended-logic payload. SPAN/tap only -- must not be
        inline.
\end{itemize}
\textbf{Rejected -- would touch the safety path (the list students must
produce):}
\begin{itemize}[leftmargin=*, nosep]
  \item \textbf{Inline IPS/DPI on the SIS conduit} -- can drop or delay a
        safety I/O packet inside the process safety time. Reject; use passive
        monitoring.
  \item \textbf{Host AV / EDR on the logic solver} -- unsupported, can consume
        scan cycles or quarantine the safety application. Reject.
  \item \textbf{Emergency patch of the certified SIS} -- invalidates the SIL
        certificate; route through safety MOC, not a security patch window.
  \item \textbf{Any control able to command the safety output} -- creates a
        spurious-trip path and a new attack surface onto the safe state.
\end{itemize}
The re-validation flag is the discriminator: only changes \emph{inside} the
SIF boundary trigger IEC~61511 re-validation; network/zoning controls outside
it do not.
\end{solutionbox}

\begin{solutionbox}{Step 8 -- IEC 61511 $\leftrightarrow$ IEC 62443 lifecycle mapping}
\vspace{1pt}\footnotesize
\renewcommand{\arraystretch}{1.2}\setlength{\tabcolsep}{4pt}
\begin{tabular}{@{}p{52mm}p{56mm}p{38mm}@{}}
\toprule
\textbf{IEC 61511 (safety)} & \textbf{IEC 62443 (security)} & \textbf{Shared gate} \\
\midrule
HAZID / HAZOP & 62443-3-2 zone\,/\,conduit + cyber risk assessment & Joint Cyber-HAZOP \\
SIF allocation (SIL-T) & SL-T allocation per zone & Risk-criteria board \\
Safety requirements spec (SRS) & Security requirements (3-3 SRs) & Combined req review \\
Design \& realisation & Secure design / hardening & Design review \\
V\&V, SIL verification & Security V\&V / pen-test & Acceptance gate \\
Operation, proof tests, MOC & Patch mgmt, monitoring, change & \textbf{Shared change board} \\
\bottomrule
\end{tabular}
\normalsize
\par\vspace{2mm}
\textbf{Load-bearing point:} any modification to the safety application is
simultaneously a safety MOC \emph{and} a security change. Neither standard's
change process may complete without the other's sign-off -- the shared change
board is the single mechanism that prevents a ``security patch'' from voiding
the SIL and a ``safety MOC'' from opening a new attack path.
\end{solutionbox}

\begin{solutionbox}{Step 9 -- (Stretch) GSN assurance-case fragment}
Acceptable indented goal tree (G\,=\,goal, S\,=\,strategy, C\,=\,context,
Sn\,=\,solution/evidence):
\begin{itemize}[leftmargin=*, nosep]
  \item \textbf{G0} HF-201 SIF-1 achieves SIL\,2 \emph{and} is not defeatable
        by a credible cyber adversary.
  \begin{itemize}[leftmargin=1.2em, nosep]
    \item \textbf{C1} Threat model: adversary on plant DMZ, TRISIS-class TTPs.
          \quad \textbf{C2} Proof-test regime $T_I=1$\,yr.
    \item \textbf{S1} Argue over safety integrity \emph{and} integrity of the
          safety function against the threat model.
    \begin{itemize}[leftmargin=1.2em, nosep]
      \item \textbf{G1} $\mathrm{PFD}_{avg}\le10^{-2}$ demonstrated.
            $\to$ \textbf{Sn1} Step-4 calculation ($2.96\times10^{-3}$).
      \item \textbf{G2} No single cyber foothold defeats $\geq2$ credited
            layers. $\to$ \textbf{Sn2} Step-7 controls (zoning, diode, signed
            download, key custody); \textbf{Sn3} Step-5 independence analysis.
      \item \textbf{G3} Integrity preserved across the lifecycle. $\to$
            \textbf{Sn4} Step-8 shared change-board gate.
    \end{itemize}
  \end{itemize}
\end{itemize}
\textbf{Marking point:} the security argument must be \emph{inside} the safety
case (G2/G3 hang under G0), not a separate annex. A SIL certificate that
ignores cyber is an incomplete assurance case -- that realisation is the
deliverable.
\end{solutionbox}

\begin{solutionbox}{Common student mistakes}
\begin{itemize}[leftmargin=*, nosep]
  \item \textbf{Treating cyber faults as random.} Plugging an attack into the
        $\beta$-factor / $\lambda_{DU}$ math. Cyber is intentional and
        correlated; it gets PFD\,$=1$ on any layer it controls, not a small
        number.
  \item \textbf{Crediting correlated layers as independent.} Missing that the
        flat LAN + shared EWS makes SIS and BPCS one common cause -- the whole
        point of Step~5.
  \item \textbf{Proposing interfering controls.} Inline IPS on the safety
        conduit, host AV on the logic solver, emergency patching the certified
        SIS. All degrade or invalidate the safety function.
  \item \textbf{Forgetting the re-validation trigger.} Not flagging that a
        change inside the SIF boundary forces IEC~61511 re-validation.
  \item \textbf{Two separate registers / cases.} Keeping safety and security
        arguments apart so the gap falls between them; GSN must integrate.
  \item \textbf{Calling the TRITON trip ``proof the SIS is safe.''} It was a
        payload defect, not a detection; a clean payload would not trip.
\end{itemize}
\end{solutionbox}

\begin{rubricbox}
Total \textbf{100 pts}; pass mark 50.
\begin{itemize}[leftmargin=*, nosep]
  \item \textbf{HAZOP rigour (20):} $\geq6$ rows, paired safety+cyber causes,
        each cyber cause correctly phase-tagged; $\geq3$ target the SIS itself.
  \item \textbf{SIL math (20):} per-subsystem $\mathrm{PFD}_{avg}$ computed,
        summed, SIL band stated, dominant term identified; within $\sim2\times$.
  \item \textbf{Independence analysis (25):} correct per-layer defeat verdict,
        names the common cause, before/after residual frequency with the
        order-of-magnitude loss. This is the core master-level skill.
  \item \textbf{Control non-interference (20):} accepted controls justified as
        non-interfering with the re-validation flag, \emph{and} an explicit
        rejected list (inline IPS, host AV, emergency patch).
  \item \textbf{Lifecycle reconciliation (15):} 61511$\leftrightarrow$62443
        mapping with the shared change board identified as the load-bearing
        gate. Bonus up to +5 for a coherent GSN fragment with the security
        argument nested inside the safety case.
\end{itemize}
\textbf{Auto-fail:} any proposed control that can delay, inhibit, or
spuriously trip the safety function and is \emph{not} flagged as rejected --
security must never compromise safety.
\end{rubricbox}

\end{document}
